\documentclass[10pt]{beamer}

% =========================
% Thème Beamer
% =========================
\usetheme{Frankfurt}        % sobre et lisible
\usecolortheme{default}

% =========================
% Packages utiles
% =========================
\usepackage[french]{babel}
\usepackage[T1]{fontenc}
\usepackage{lmodern}


\RequirePackage{tabularray}
\RequirePackage{xcolor}

% Supprimer "Table X:" dans les légendes tabularray
\DefTblrTemplate{caption-tag}{default}{}
\DefTblrTemplate{caption-sep}{default}{}
\DefTblrTemplate{caption-text}{default}{}

\usepackage{wasysym}

\newcommand{\ok}[1]{\textcolor{teal}{#1}}
\newcommand{\nok}[1]{\textcolor{red}{#1}}
\newcommand\mydots{\makebox[1em][c]{.\hfil.\hfil.}}

% =========================
% Infos du document
% =========================
\title[Cas pratiques]{Séminaire d'application PCT}
\author{Aurélien Gabieli}
\date{\today}

% =========================
% Footer : numéro de slide
% =========================
\setbeamertemplate{footline}[frame number]


\usepackage{hyperref}

\newcommand{\pctart}[1]{%
  \href{https://www.wipo.int/pct/fr/texts/articles/a#1.html}{Art.~#1}%
}

\newcommand{\pctrule}[1]{%
  \href{https://www.wipo.int/pct/fr/texts/rules/r#1.html}{Règle~#1}%
}

\hypersetup{
  %colorlinks=true,
  %linkcolor=white,
  urlcolor=blue
}


% =========================
% Sommaire automatique
% =========================

% Slide plan au début
\begin{document}

\begin{frame}
  \titlepage
\end{frame}

\begin{frame}{Plan}
  \tableofcontents
\end{frame}

% =========================
% Plan au début de chaque section
% =========================
\AtBeginSection[]{
  \begin{frame}{Plan}
    \tableofcontents[currentsection]
  \end{frame}
}

% ======================
% Examens
% ======================

% =========================
% Examen 2023
% =========================

\section[Examen 2023]{Examen 2023}

\begin{frame}{}

\begin{block}{}

\end{block}

\begin{itemize}
    \item En l'espèce, 
\end{itemize}
\end{frame}
% =========================
% Examen 2024
% =========================

\section[Examen 2024]{Examen 2024}

%--------------------------
\begin{frame}{Question 1}

\begin{block}{}
La demande internationale cesse de produire ses effets en phase nationale si :

\begin{itemize}
  \item[$\circ$] a. le délai de 30 mois à compter de la date de priorité n’est pas respecté pour 
  l’ouverture de la phase nationale. 
  \item[$\circ$] b. la taxe nationale n’est pas payée dans le délai de 31 mois à compter de la date de 
  priorité. 
  \item[$\circ$] c. les signatures manquantes ne sont pas remises dans le délai de 30 ou 31 mois à 
  compter de la date de priorité. 
  \item[$\bullet$] d. Aucune des réponses a, b ou c n’est correcte.
\end{itemize}

\end{block}

\end{frame}

%--------------------------
\begin{frame}{Question 1}

\begin{itemize}
  \item[$\circ$] \textbf{a.} \; Faux : le délai d’entrée en phase nationale n’est pas toujours
de 30 mois, certains États, LU et TZ, appliquant encore un délai de 20 mois
(\pctart{22}.1)).

  \item[$\circ$] \textbf{b.} \; Faux : certains États, LU et TZ, appliquant encore un délai de 20 mois
  (\pctart{22}.1)).

  \item[$\circ$] \textbf{c.} \; Faux : la signature de tout
  déposant qui n'a pas signé la requête est une condition
  permise par le PCT mais pas obligatoire
  (\pctart{27}.2) + \pctrule{51bis}.1.a.i)) 

  \item[$\circ$] \textbf{d.} \; Correct puisqu'aucune des réponses \textbf{a},
  \textbf{b} ou \textbf{c} est correcte.
\end{itemize}

\end{frame}

%--------------------------
\begin{frame}{Question 2}

\begin{block}{}
L’absence de dépôt d’une demande d’examen préliminaire international :

\begin{itemize}
  \item[$\bullet$] a. n’a aucun effet sur le délai d’ouverture de la phase nationale. 
  \item[$\bullet$] b. ne ralentit pas nécessairement les procédures nationales de délivrance en phase 
  nationale. 
  \item[$\bullet$] c. empêche le déposant de déposer des modifications de la description ou des dessins 
  pour qu’elles soient examinées pendant la phase internationale. 
  \item[$\circ$] d. Aucune des réponses a, b ou c n’est correcte. 
\end{itemize}

\end{block}

\end{frame}

%--------------------------
\begin{frame}{Question 2}

\textbf{Base juridique :} \pctart{22}(1) + \pctart{34}(2)(b).  

\begin{itemize}
  \item[$\bullet$] \textbf{a.} \; Correct : le délai d’entrée en phase nationale
  n’est pas affecté par l’absence de Chapitre II (\pctart{22}(1)).

  \item[$\bullet$] \textbf{b.} \; Correct : l’absence d’examen préliminaire n’implique pas nécessairement
  un ralentissement, car l’examen reste national.

  \item[$\bullet$] \textbf{c.} \; Correct : les modifications de la description et des dessins pendant la phase
  internationale sont examinées uniquement dans le cadre du Chapitre II
  (\pctart{34}(2)(b)).

  \item[$\circ$] \textbf{d.} \; Faux puisque \textbf{a}, \textbf{b} et \textbf{c} sont correctes.
\end{itemize}

\end{frame}

%--------------------------
\begin{frame}{Question 3}

\begin{block}{}
L’office récepteur est seul compétent pour décider :

\begin{itemize}
  \item[$\bullet$] a. d’exiger une taxe de transmission pour les demandes PCT déposées auprès de lui. 
  \item[$\bullet$] b. d’accepter une seule langue ou plusieurs langues dans laquelle ou dans lesquelles 
  des demandes PCT peuvent être déposées auprès de lui. 
  \item[$\bullet$] c. du choix des administrations chargées de la recherche internationale supplémentaire 
  auxquelles les déposants auront accès s’ils déposent auprès de lui. 
  \item[$\circ$] d. Aucune des réponses a, b ou c n’est correcte. 
\end{itemize}

\end{block}

\end{frame}

%--------------------------
\begin{frame}{Question 3}

\textbf{Base juridique :} \pctrule{14}.1(c) + \pctrule{12}.1(a) + \pctrule{45bis}.9.  

\begin{itemize}
  \item[$\bullet$] \textbf{a.} \; Correct : chaque office récepteur décide d’exiger ou non
  une taxe de transmission (\pctrule{14}.1(c)).

  \item[$\bullet$] \textbf{b.} \; Correct : la langue de dépôt est déterminée par ce que
  l’office récepteur accepte (\pctrule{12}.1(a)).

  \item[$\bullet$] \textbf{c.} \; Correct : l’accès aux administrations chargées de la recherche
  supplémentaire dépend de l’office récepteur compétent (\pctrule{45bis}.9).

  \item[$\circ$] \textbf{d.} \; Faux puisque \textbf{a}, \textbf{b} et \textbf{c} sont correctes.
\end{itemize}

\end{frame}

%--------------------------
\begin{frame}{Question 4}

\begin{block}{}
Sous réserve que les autres exigences de l’article 11(1) soient respectées,
l’office récepteur attribuera une date de dépôt international à condition que :

\begin{itemize}
  \item[$\bullet$] a. au moins un déposant soit domicilié dans un État contractant du PCT. 
  \item[$\bullet$] b. au moins un déposant soit domicilié dans un État contractant du PCT ou soit 
  ressortissant d’un État contractant du PCT. 
  \item[$\bullet$] c. au moins un déposant soit domicilié dans un État contractant du PCT et soit 
  ressortissant d’un État contractant du PCT. 
  \item[$\circ$] d. Aucune des réponses a, b ou c n’est correcte. 
\end{itemize}

\end{block}

\end{frame}

%--------------------------
\begin{frame}{Question 4}

\begin{itemize}
  \item[$\circ$] \textbf{a.} \; Correct.

  \item[$\bullet$] \textbf{b.} \; Correct : au moins un déposant doit avoir
  nationalité \emph{ou} domicile dans un État contractant
  (\pctart{11}(1)(i) + \pctart{9}(1)).

  \item[$\circ$] \textbf{c.} \; Correct.

  \item[$\circ$] \textbf{d.} \; Faux.
\end{itemize}

\end{frame}

%--------------------------
\begin{frame}{Question 5}

\begin{block}{}
La date de dépôt d’une demande nationale antérieure sera reconnue aux fins de la 
phase internationale comme date de priorité dans une demande PCT :

\begin{itemize}
  \item[$\bullet$] a. si la priorité de ladite demande antérieure est valablement revendiquée dans la 
  demande PCT et qu’elle respecte les exigences de la Convention de Paris. 
  \item[$\circ$] b. même si la date de dépôt de cette demande nationale est antérieure de plus de 12 
  mois à la date de dépôt PCT, et qu’elle respecte les autres exigences de la 
  Convention de Paris. 
  \item[$\bullet$] c. même si la date de dépôt de cette demande nationale est antérieure de plus de 12 
  mois à la date du dépôt PCT, mais à condition qu’elle tombe dans le délai de 12 + 2 
  mois avant ladite date de dépôt PCT, et qu’elle respecte les autres exigences de la 
  Convention de Paris. 
  \item[$\circ$] d. Aucune des réponses a, b ou c n’est correcte. 
\end{itemize}

\end{block}

\end{frame}

% %--------------------------
% \begin{frame}{Question 5 - à revoir}

% \begin{itemize}
%   \item[$\bullet$] \textbf{a.} \; Correct : une priorité valablement revendiquée dans les 12 mois
%   produit effet conformément à la Convention de Paris
%   (\pctart{8}(1) + \pctart{8}(2)(a)).

%   \item[$\circ$] \textbf{b.} \; Faux : une priorité revendiquée après l’expiration des 12 mois
%   n’est pas reconnue automatiquement.

%   \item[$\bullet$] \textbf{c.} \; Correct : une restauration du droit de priorité est possible
%   si la demande PCT est déposée dans un délai de 2 mois après l’expiration des 12 mois
%   (\pctrule{26bis}.3).

%   \item[$\circ$] \textbf{d.} \; Faux puisque \textbf{a} et \textbf{c} sont correctes.
% \end{itemize}

% \end{frame}

%--------------------------
\begin{frame}{Question 6}

\begin{block}{}
La question de l’unité de l’invention :

\begin{itemize}
  \item[$\bullet$] a. n’est pas examinée par l’office récepteur lors de la vérification des taxes acquittées. 
  \item[$\bullet$] b. n’est pas examinée par l’administration chargée de la recherche internationale après 
  l’établissement du rapport de recherche internationale. 
  \item[$\bullet$] c. peut être examinée par l’administration chargée de l’examen préliminaire 
  international même si l’administration chargée de la recherche internationale ne l’a 
  pas examinée. 
  \item[$\circ$] d. Aucune des réponses a, b ou c n’est correcte. 
\end{itemize}

\end{block}

\end{frame}

%--------------------------
\begin{frame}{Question 6}

\begin{itemize}
  \item[$\bullet$] \textbf{a.} \; Correct : l’office récepteur vérifie les taxes formelles
  mais ne procède pas à un examen de l’unité de l’invention (\pctrule{13}).

  \item[$\bullet$] \textbf{b.} \; Correct : l’examen de l’unité par l’ISA intervient pendant la recherche,
  pas après l’établissement du rapport de recherche internationale
  (\pctart{17}.3.a)).

  \item[$\bullet$] \textbf{c.} \; Correct : l’IPEA peut examiner l’unité indépendamment
  de l’ISA et inviter au paiement de taxes additionnelles
  (\pctart{34}.3.a)).

  \item[$\circ$] \textbf{d.} \; Faux puisque \textbf{a}, \textbf{b} et \textbf{c} sont correctes.
\end{itemize}

\end{frame}

% =========================
% Cas pratique PCT
% =========================

\subsection{Cas pratique}

%--------------------------
\begin{frame}{Cas pratique a)}

\begin{block}{quand?}
Délai de priorité : 12 mois (\pctart{8}.2.a)) + Art. 4.C.1 CUP.
\end{block}

En l'espèce, une demande nationale a été déposée le 17 mai 2023.\newline

\textbf{Conclusion: dépôt PCT sous priorité avant le 17/05/2024.}

\end{frame}

%--------------------------
\begin{frame}{Cas pratique a)}

\begin{block}{où?}
La demande internationale doit être déposée auprès d’un \textbf{office récepteur compétent},
c’est-à-dire :

\begin{itemize}
  \item un office national de l’État PCT où le demandeur est \textbf{domicilié}
  (\pctrule{19}.1(a)(i)) ;

  \item un office national de l’État PCT dont le demandeur est \textbf{ressortissant}
  (\pctrule{19}.1(a)(ii)) ;

  \item le Bureau international agissant en tant qu’office récepteur (RO/IB)
  (\pctrule{19}.1(a)(iii)) ;

  \item une organisation intergouvernementale (ex. OEB, ARIPO) ou un office national
  agissant pour l’État où un demandeur est domicilié ou ressortissant
  (\pctrule{19}.1(b)).
\end{itemize}
\end{block}

En l'espèce, A-FR est une société française. \newline

\textbf{Conclusion : RO = INPI, OEB, IB}

NB: pas d'obligation de dépôt RO/FR car pas première demande.

\end{frame}

\begin{frame}{Cas pratique a)}

\begin{block}{comment?}
\begin{itemize}
    \item déclaration de priorité: \pctrule{4}.10.a) + \pctrule{4}.1.b.i) (formulaire)
    \item demande internationale: \pctart{3}.2 (DI) + \pctart{4}.1 (requête) + \pctart{4}.1.d) (signatures).
\end{itemize}
\end{block}

\textbf{Conclusion: déposer une demande internationale comprenant:}
\begin{itemize}
    \item une requête comprenant une pétition PCT, désignations,
    nom \& renseignements du déposant et mandataire, titre invention,
    nom \& renseignements de l'inventeur;
    \item la requête est signée par déposant et mandataire;
    \item la copie de la demande FR.
\end{itemize}
NB: De préférence, requête sur formulaire PCT/RO/101.
\end{frame}

%--------------------------
\begin{frame}{Cas pratique b)}

\begin{block}{}
\begin{itemize}
  \item ajouter déposant après dépôt : \pctrule{92bis}.
  \item publication anticipée \pctart{21}.2.b).
\end{itemize}
\end{block}

En l'espèce, B-IE souhaite être rendue publique dès que possible.\newline

\textbf{Conclusion: ajouter co-déposant B-IE et
avancer la publication.}\newline
NB: éventuellement payer une taxe spéciale
si RRI pas disponible (\pctrule{48}.4) mais budget raisonnable.
\end{frame}

%--------------------------
\begin{frame}{Cas pratique c)}

\begin{block}{}
entrée en phrase nationale: \pctart{22}
\end{block}

En l'espèce, un produit suspect apparaît aux Pays-Bas.
Le niveau de développement de l'état de l'art est
nettement inférieur donc invention brevetable.
Pas de comparaison sur l'objet litigieux, éviter
une modification des revendications.\newline

\textbf{Conclusion : entrée en phase anticipée DO/EP.}\newline
NB: Etape suivante logique: valider brevet EP en EP/NL.
\end{frame}

%--------------------------
\begin{frame}{Cas pratique d)}

\begin{block}{}
\begin{itemize}
    \item entrée en phase anticipée \pctart{23}.3 + \pctart{40}.2
    \item PPH possible après ISR (pratique).
\end{itemize}
\end{block}

En l'espèce, CA-P veut accélérer la délivrance en ZA, US et MA.\newline

\textbf{Conclusion: ISR par EP => entrée en phase PPH(IP5) CN et US => AP Accord bilatéral CN-AP.}\newline
NB: PCT4ever

\end{frame}

% =========================
% Examen 2025
% =========================

\section[Examen 2023]{Examen 2025}

\begin{frame}{}

\begin{block}{}

\end{block}

\begin{itemize}
    \item En l'espèce, 
\end{itemize}
\end{frame}

\end{document}